%% Annexe A : Glossaire
\chapter{Glossaire}
\label{ann:glossaire}

Ce glossaire reprend les notions définies dans les chapitres. Le renvoi
indique le chapitre où la notion est présentée.

\begin{description}
    \item[Actif] Élément de valeur pour l'organisation, qu'il faut protéger (Ch.~\ref{chap:cia}).
    \item[Appétence au risque] Niveau de risque qu'une organisation est prête à accepter (Ch.~\ref{chap:vocabrisque}).
    \item[Appréciation des risques] Démarche d'identification, d'analyse et d'évaluation des risques (Ch.~\ref{chap:methodes}).
    \item[Déclaration d'applicabilité (DdA)] Document indiquant, pour chaque mesure de l'Annexe~A, si elle est retenue, sa justification et son état (Ch.~\ref{chap:dda}).
    \item[Disponibilité, Intégrité, Confidentialité (DIC)] Les trois critères fondamentaux de la sécurité de l'information (Ch.~\ref{chap:cia}).
    \item[Gouvernance de la sécurité] Responsabilités et pratiques par lesquelles la direction oriente la sécurité (Ch.~\ref{chap:gouvernance}).
    \item[Mesure de sécurité] Moyen (organisationnel, humain, physique, technique) de traiter un risque (Ch.~\ref{chap:iso27002}).
    \item[Non-conformité] Écart entre une exigence et la réalité, majeure ou mineure (Ch.~\ref{chap:audit}).
    \item[Partie intéressée] Personne ou organisme ayant une attente vis-à-vis du SMSI (Ch.~\ref{chap:clauses46}).
    \item[PCA / PRA] Plan de continuité d'activité~; plan de reprise d'activité (Ch.~\ref{chap:continuite}).
    \item[PDCA] Cycle Planifier-Déployer-Contrôler-Agir, moteur de l'amélioration continue (Ch.~\ref{chap:pdca}).
    \item[Politique de sécurité (PSSI)] Document fondateur fixant objectifs, règles et responsabilités (Ch.~\ref{chap:politique}).
    \item[Revue de direction] Examen périodique du SMSI par la direction, qui décide des améliorations (Ch.~\ref{chap:audit}).
    \item[Risque] Combinaison de la probabilité d'un événement et de son impact (Ch.~\ref{chap:vocabrisque}).
    \item[Risque résiduel] Risque subsistant après traitement, accepté par la direction (Ch.~\ref{chap:traitement}).
    \item[RTO / RPO] Délai de reprise visé~; perte de données maximale acceptable (Ch.~\ref{chap:continuite}).
    \item[SMSI] Système de management de la sécurité de l'information (Ch.~\ref{chap:prologue}).
    \item[Vraisemblance] Probabilité qu'un risque se réalise (Ch.~\ref{chap:vocabrisque}).
\end{description}
